This paper examines whether artificial curiosity constitutes a genuine analog of human curiosity or a sophisticated simulation of it. We review the neuroscience of human curiosity, including information gap theory \citep{loewenstein1994}, the Five-Dimensional Curiosity Scale Revised (5DCR) framework \citep{kashdan2018}, and developmental research demonstrating curiosity's embodied, multimodal character. We then survey AI curiosity mechanisms across three paradigms: intrinsic motivation in reinforcement learning (prediction-error, count-based, and random network distillation approaches), behavioral evaluation of curiosity in large language models, and developmental robotics work employing active inference. Cross-cultural evaluation of LLM curiosity reveals a 63\% compression of human curiosity diversity, with models defaulting to Western-normative, socially neutral question patterns \citep{borah2025}. However, curiosity-driven developmental robotics produces convergent developmental dynamics---including U-shaped regression curves and rote-before-compositional learning stages---that mirror human child development without explicit curriculum design \citep{tinker2025}. We conclude that current theoretical tools cannot distinguish genuine from simulated curiosity, as the distinction requires a theory of phenomenal consciousness that does not yet exist. The functional convergence observed in developmental dynamics suggests the boundary between genuine curiosity and its simulation may be less sharp than commonly assumed.
